\section{Response to Reviewer 1}
 {\color{blue}
  We sincerely thank you for your insightful review of our manuscript and especially for your detailed, constructive, and insightful comments on the organization of the proofs.
  We have revised the manuscript accordingly and have answered all your questions in detail.

  In particular, we have improved the organization and notations of the paper, clarified the algorithmic descriptions, added missing proof details, corrected the typos and derivation issues you pointed out.  We have also improved the readability of the proofs. We believe that the presentation is now in good shape to convey the main ideas of the paper.

  We have also provided additional explanations for the numerical experiments. We are grateful for your suggestions, which helped us substantially improve the clarity, rigor, and presentation of the paper.

 }
\subsection{Major concerns}

\paragraph{General presentation of the paper}
\begin{itemize}
      \item \comm{Definition 2.1 can be introduced earlier, as the authors refer to that in the introduction, for example on Page 1, line 55.}

            \resp{Thank you for pointing this out. We have revised the introduction to avoid referring to Definition~2.1 before it is introduced.}

      \item \comm{Assumption 2.3: the authors can comment that this is equivalent to Lipschitz gradient.}

            \resp{Thank you. We have added this comment.}

      \item \comm{Notation: please include more details about the notation used in the Preliminaries section, on Page 4, line 4. For example, the authors denote by $x_k$, $v_k$, and $s_k$ the iterates of Algorithm 1. These quantities are used in the lemmas, for example Lemma 2.3, and it is not obvious whether they refer to the iterates from the algorithm or to general vectors.}

            \resp{Thank you. We have added the suggested details to clarify that $x_k$, $v_k$, and $s_k$ refer to the iterates generated by Algorithm~1.}

      \item \comm{The authors sometimes use the same letter for different notation across the paper. I suggest avoiding this. For example, $\lambda_k$ is used as the dual variable for the constraint $\|d_k\| \le r_k$ in TR+, while in Lemma 3.6 the notation $\operatorname{diag}(\lambda_1,\ldots,\lambda_n)$ is used to represent eigenvalues.}

            \resp{Thank you. We have reviewed the notation throughout the paper and revised it to ensure consistency.}

      \item \comm{Please add text indicating that the proof of Lemma 4.4 appears in Appendix C.}

            \resp{Thank you. We have added the corresponding text.}

      \item \comm{Please add text indicating that the proof of Lemma 4.7 appears in Appendix C.}

            \resp{Thank you. We have added the corresponding text.}
\end{itemize}

\paragraph{Proofs and mathematical derivations}

\comm{In all proofs, the authors are encouraged to provide more detailed explanations of the transitions from one step to another, even if some of those transitions are trivial.}

\resp{Thank you for your comments. We carefully reviewed all of our proofs and revised the manuscript accordingly. To improve the flow, the analysis for each variant is now organized into iteration complexity, oracle calls of the subroutine, and finally, oracle complexity. We add more details to the proofs to improve the clarity. \\
      For comments where the corresponding revisions are straightforward and self-explanatory, we have omitted detailed responses for brevity and adopted proper revision in the manuscript. Regarding your comment on Lemma 3.13, we have completely rewritten the result (now Lemma 3.11) and improved its dependence on $\nu$.}
\begin{itemize}
      \item \comm{Page 7, line 16. Even if it is obvious, mention when introducing the sequence $\{\phi_k(x)\}_{k\ge 0}$ that it is a monotone non-increasing sequence, since this property is used in Lemma 3.4.}

            \resp{Thanks for the comment. We do not use, nor claim, that $\phi_k(x)$ is non-increasing; this is generally false. The proof only relies on the way we define the estimating sequence.}

      \item \comm{Lemma 3.4, Page 8, line 7. Please provide more details about this equality. The manuscript says that ``the fifth line is from rearranging terms and solving the trivial minimization problem,'' but it would be better to provide more details, for example in the appendix.}

            \resp{We have added all the details in this proof.}

      \item \comm{From the statement of Lemma 3.5, by combining the bound on $\|d_+\|$ and the bound from $\|\nabla f(y_k)\|$, we obtain $\|d_+\|\le \kappa_H/M$. Does this affect any of the results for $k\in\mathcal{L}$?}

            \resp{Thanks for the detailed observation. No, it doesn't contradict with any results.}

      \item \comm{Lemma 4.1, Page 15, line 42. For the last inequality, please add more details about the derivation. For example, mention that $\|d_-\|\le r_-$, where $r_-=\frac{1}{M}\sigma_-$ from Subroutine 2, line 3. The same applies to the value of $\sigma_-$. Also, in this lemma, the proof uses $\sigma_-\le \sigma_-^2$ in the last inequality, but that is not clarified. From the value of $\sigma_-$ in Subroutine 2, line 3, this is not clear unless the assumption $2M\epsilon\ge 1+2\theta$ is made.}

            \resp{Thanks for the question. We did not use $\sigma_-\leq \sigma_-^2$, and we have added clarification in the revision.}

      \item \comm{Proof of Lemma 2.3. The proof implicitly assumes $M\ge 1$, but this is not mentioned in the paper. For example, Assumption 2.1 assumes $M>0$, but for the proof to work, it needs $M\ge 1$. The second inequality is in fact based on $\mu\ge \sqrt{M}\|d\|$, rather than $\mu\ge M\|d\|$. The same applies to line 28. It should be $\frac{3}{4}M\|d\|^4$ instead of $\frac{3}{4}M^2\|d\|^4$, unless one assumes $1\le M\le M^2$ and hence $-\frac{M}{4}\|d\|^4\ge -\frac{M^2}{4}\|d\|^4$. This needs to be fixed and it affects other lemmas, including Lemma 3.3 and others. The fix is to assume $M\ge 1$, so one can use $\mu\ge M\|d\|\ge \sqrt{M}\|d\|$ for the first part, and for the second part use
                  $-\frac{M}{4}\|d\|^4\ge -\frac{M^2}{4}\|d\|^4$.}

            \resp{Thank you for looking into this. We did not use $M>1$. The second inequality is based on that
                  \begin{equation*}
                        M^2\|d\|^2 \leq \mu^2, \text{and thus } M^2\|d\|^4\leq \mu^2\|d\|^2.
                  \end{equation*}
                  The same applies to line 28
                  \begin{equation*}
                        4M\|d\| \langle \nabla f(x+d),-d \rangle \geq_{(a)} 2\mu \langle \nabla f(x+d),-d \rangle \geq \|\nabla f(x+d)\|^2 +\frac{3}{4}\mu^2 \|d\|^2\geq_{(b)}  \|\nabla f(x+d)\|^2 +\frac{3}{4} M^2 \|d\|^4,
                  \end{equation*}
                  where $(a)$ is from $\mu \leq 2M\|d\|$, $(b)$ is from $\mu \geq M\|d\|$.
            }

      \item \comm{Proof of Lemma 3.9. There are issues with the second and third inequalities. For example, based on Lemma B.1, the second inequality should involve
                  $\|(\nabla^2 f(y)+\xi I)^{-1}\nabla f(y)\|^3$.
                  Also, the transition from
                  $\sum_{i=1}^n \frac{\beta_i}{(\lambda_i+\xi)^3}$
                  to
                  $\sum_{i=1}^n \frac{\beta_i}{(\lambda_i+\xi)^2}$
                  in order to replace
                  $\sum_{i=1}^n \frac{\beta_i}{(\lambda_i+\xi)^2}$
                  by $\|(\nabla^2 f(y)+\xi I)^{-1}\nabla f(y)\|$ is not properly done.}

            \resp{We have now rewrote the proof accrording to your comment.}

      \item \comm{Lemma B.2 and its proof, Page 24, line 14. The definition $\eta_0=Ca_0$ is not the same as the definition of $\eta_0$ under Corollary 3.2, equation (3.5). Please clarify.}

            \resp{We have rewrote the proof.}

      \item \comm{Lemma B.2 and its proof, Page 24, line 30. First, it is not clear how $w\le w^{-2}$ gives $\ln(1/w)\ge 2$. Also, in line 30, the proof uses $\ln(1/w)\le 2$ for the inequality to hold.}

            \resp{Thank you for pointing this out. We have refined our proof.}

      \item \comm{Lemma C.3, Page 28, line 5. There is a minor calculation mistake in the computation of the derivative. It should be $4A$ instead of $2A$. This affects the final bound in the lemma.}

            \resp{Thanks. We have revised it and the corresponding bound.}

      \item \comm{Lemma C.5. There is a typo in the definition of $C(G_0,M_0,M,\epsilon)$; I think there is an extra $M$. Please also recompute the bound on lines 10--12. I think there is a minor computation mistake: the first term should be $(1+2\theta)^3$ instead of $(1+2\theta)^{3/2}$, and $\eta^2$ instead of $\eta$; also, for the fourth term, it should be $(1+2\theta)^2$ instead of $(1+2\theta)$, and there should be $\eta^{3/2}$ in the denominator.}

            \resp{Thanks. We have recomputed the bound as you suggested and added more details to the proof.}
\end{itemize}

\paragraph{Algorithm}
\begin{itemize}
      \item \comm{Subroutine 1, Local Detection (LD), lines 2--4. Can you link it to the results in Corollary 3.1 when describing the algorithm on Page 6, lines 40--53, where $\lambda_k=0$ triggers the LD subroutine?}

            \resp{We now added some discussion in the description of the algorithm.}

      \item \comm{Subroutine 1, Local Detection (LD), line 20. Can you comment on why TR+ is called with $\sigma=0$, i.e., $\operatorname{TR}^+(y,0,r)$?}

            \resp{This is because acceleration requires $\frac{\mu}{\|d\|} = \Theta(1)$. At this stage, we have already identified the bracketing points for the bisection procedure, so the remaining task is to adjust the radius. Therefore, it suffices to call $\operatorname{TR}^+(y,0,r)$ with $\sigma=0$.}

      \item \comm{Subroutine 2, Ratio Bracketing and Bisection (R\&B) for Algorithm 2. There are typos in line 5: $y$ and $d$ are not defined. They should probably be $y(\sigma_-)$ and $d_-$.}

            \resp{Thanks. We have corrected it.}

      \item \comm{Subroutine 2. There are typos: the condition $\theta>1$ should be the same as in Algorithm 2. The same applies to $0<\eta<1$.}

            \resp{Thanks. We have fixed it.}
\end{itemize}

\paragraph{Numerical Experiments}
\begin{itemize}
      \item \comm{In the Contribution section, Page 3, line 28, the paper refers to [42] and states: ``On a global scale, all accelerated SOM, including ours and [42, 46].'' However, the plots do not show the performance of running [42]. Please either remove that statement or add the results of running this method.}

            \resp{Thanks. After careful consideration, we have narrowed the reference to [46]. The original MS method (A-NPE) in [42], together with the searching procedure in Section 7 of [42], is primarily intended as a general algorithmic framework rather than a practical implementation, and its practical limitations have also been discussed in Section 4.3.3 of [47]. In practice, competitive MS-type methods require substantial enhancements \citep{carmon2022optimal}. Since our primary objective is to investigate the acceleration of trust-region-type methods and because the space limitation here, we have chosen to focus on the experiments of our approaches accordingly. A systematic empirical study of MS-type methods with different second-order oracles, including the $\trp$ oracle, is an interesting direction that we plan to pursue in future work.
            }

      \item \comm{Please provide a link to the Git source code. This will help others use the method and compare against it.}

            \resp{Thanks for the suggestion. We have added a link to our current implementation under long-term maintenance: \url{https://github.com/bzhangcw/UTR.jl}.}

      \item \comm{In the plots, the paper uses ATR and ATR (MS), but these two notations are never introduced. I think ATR and ATR (MS) refer to Algorithm 1 and Algorithm 2, but it would be easier to define these, especially for readers who skim the paper and check the plots directly.}

            \resp{Thanks for this suggestion. We have changed the legend names to match the algorithm names.}

      \item \comm{I recommend adding one line to quickly explain the difference between the datasets, for example the size of each one, $N$, or the ratio of both classes. The motivation behind running all algorithms on these different datasets is currently unclear.}

            \resp{Thanks for this suggestion. We add a sentence to explain the difference between the datasets. Each row of presented figures share the sparsity and dimension but different sample size.}

      \item \comm{Please provide extra details on the $\epsilon$ value used and the $\theta$ and $\eta$ parameter values used for Algorithm 2.}

            \resp{Thanks for this suggestion. In the implementation, $\eta = 0.3$, $\theta = 1.3$, and $\epsilon_g$ is the precision parameter. Details can be found in the git repo.}

      \item \comm{Please provide extra details on the computation of $D_0$ and $G_0$ in the implementation of Algorithm 2, since these quantities are required for Subroutine 2. For example, how did you compute $x^\ast$?}
            \resp{Thank you for this suggestion. In our implementation, we do not compute $x^*$ and therefore do not explicitly evaluate $D_0=\|x_0-x^*\|$ or the corresponding theoretical bound $G_0$. These quantities are used only in the analysis to certify the existence of a valid upper bracketing point $\sigma_+$. In practice, we start from a prescribed positive guess for $\sigma_+$, solve the trust-region subproblem at $y(\sigma_+)$, and check the verifiable upper-bracketing condition
                  \[
                        \lambda=0,
                        \qquad
                        \|d\|\leq\frac{\eta}{M}\sigma_+.
                  \]
                  If this condition is not satisfied, we increase $\sigma_+$ by a fixed positive increment and repeat the test until it holds. Thus, the numerical implementation requires no knowledge of $x^*$, $D_0$, or $G_0$. This adaptive bracketing idea is analogous to the search procedure of \citet[Section~7]{monteiro2013accelerated}.}
\end{itemize}

\subsection{Minor concerns}
\begin{itemize}
      \item \comm{I think you can mention in the abstract that this work is only for unconstrained convex optimization.}

            \resp{Thanks. We now mention in the abstract that this work focuses on unconstrained convex optimization.}

      \item \comm{Page 1, line 55. ``needed to find an $\epsilon$-approximate solution.'' $\epsilon$ was not defined yet, and it would be better to add a sentence mentioning that $\epsilon$ represents the gradient termination tolerance, for example.}

            \resp{Thanks. We now removed the notion ``$\epsilon$-approximate'' in the introduction.}

      \item \comm{Lemma 2.3, Page 4, line 50. Typo: ``hold, holds'' should be ``hold''.}

            \resp{Thanks. We have corrected this typo.}

      \item \comm{Lemma 3.14, Page 13, line 21. Typo: ``holds'' should be ``hold''.}

            \resp{Thanks. We have corrected this typo.}

      \item \comm{In Section 4, please replace $a_k$ with $a_k(\sigma_k)$, since this is how $a_k(\sigma_k)$ is defined in Section 4.2, Page 14, line 19. The same applies to $y_k(\sigma_k)$ instead of $y_k$, as defined in (4.1).}

            \resp{Thanks. We have replaced $y_k$ with $y_k(\sigma_k)$ in Section 4 and add definition of $a_k$ at line 7 of Algorithm 2.}

      \item \comm{Proof of Lemma 3.9. For clarity, please move the definition of $\xi\in[\mu_l,\mu_u]$ to line 45 instead of line 46. For example, one can write: Using $\xi\in[\mu_l,\mu_u]$ and applying the mean value theorem to $r_y(\mu_l)-r_y(\mu_u)$, we have \ldots}

            \resp{Thanks. We have revised the proof as suggested.}

      \item \comm{Proof of Lemma 4.3, Page 24, line 43. Please refer to Algorithm 2, line 9.}

            \resp{Thanks. We have added the reference to Algorithm 2, line 9.}

      \item \comm{Proof of Lemma 4.3, Page 25, line 11. Typo: there should be a factor of $6$ when taking the derivative of the Lagrangian. This factor cancels out at the end, and the final result holds.}

            \resp{Thanks. We have corrected the derivative by adding the missing factor of $6$.}

      \item \comm{Lemma C.2. There is a typo in (C.4). It should be $\delta$ instead of $r$.}

            \resp{Thanks. We have corrected this typo.}

      \item \comm{Proof of Lemma 4.7. Please add ``using (2.4)'' to line 49.}

            \resp{Thanks. We have added ``using (2.4)'' as suggested.}
\end{itemize}
